\documentclass[a4paper,12pt]{article}

\usepackage[T2A]{fontenc}
\usepackage[utf8]{inputenc}
\usepackage[russian]{babel}
\usepackage{amsmath,amssymb,amsthm,mathtools}
\usepackage{helvet}
\renewcommand{\familydefault}{\sfdefault}
\usepackage{icomma}
\usepackage[left=18mm,right=18mm,top=18mm,bottom=20mm]{geometry}
\usepackage{microtype}
\usepackage{tikz}
\usepackage{tkz-euclide}
\usepackage{pgfplots}
\pgfplotsset{compat=1.18}
\usepackage{tabularx,booktabs,longtable}
\usepackage{enumitem}
\usepackage[hidelinks]{hyperref}
\usepackage{parskip}
\usepackage{fancyhdr}
\usepackage{setspace}
\usepackage{xcolor}

\definecolor{accent}{HTML}{008080}
\definecolor{darkaccent}{HTML}{004D4D}
\definecolor{textgray}{HTML}{333333}
\definecolor{softgray}{HTML}{F3F6F6}

\providecommand{\tg}{\operatorname{tg}}
\providecommand{\ctg}{\operatorname{ctg}}
\providecommand{\arctg}{\operatorname{arctg}}

\newcommand{\Z}{\mathbb{Z}}
\newcommand{\important}[1]{\textcolor{darkaccent}{\textbf{#1}}}

\pagestyle{fancy}
\fancyhf{}
\lhead{\textcolor{textgray}{Тригонометрические уравнения}}
\rhead{\textcolor{textgray}{ЕГЭ}}
\cfoot{\thepage}
\renewcommand{\headrulewidth}{0.4pt}
\renewcommand{\headrule}{\hbox to\headwidth{\color{accent}\leaders\hrule height \headrulewidth\hfill}}

\onehalfspacing
\setlist[itemize]{leftmargin=1.4em,itemsep=0.25em}
\setlist[enumerate]{leftmargin=1.6em,itemsep=0.35em}

\begin{document}

\begin{titlepage}
\thispagestyle{empty}
\centering
\vspace*{0.18\textheight}

{\Huge\bfseries\textcolor{darkaccent}{Чек-лист}\par}
\vspace{0.8em}
{\LARGE Тригонометрические преобразования и уравнения\par}

\vspace{2.5em}
{\large Лёвин Артём Александрович эксклюзивно для Тимофея Васильченко\par}

\vspace{1.2em}
{\large \today\par}

\vfill

\begin{tikzpicture}[x=1cm,y=1cm]
\draw[accent,line width=1.2pt]
  (-5.5,0) .. controls (-3.0,0.85) and (-1.4,-0.75) .. (0.8,0.1)
  .. controls (2.9,0.9) and (4.0,-0.5) .. (5.5,0.18);
\draw[darkaccent,line width=0.4pt]
  (-5.5,-0.25) .. controls (-2.5,0.45) and (1.2,-0.8) .. (5.5,-0.05);
\end{tikzpicture}
\end{titlepage}

\section*{1. Главная идея темы}

Тригонометрическое уравнение обычно решается через предварительное упрощение:
\[
\text{формулы} \quad \longrightarrow \quad
\text{одинаковые углы} \quad \longrightarrow \quad
\text{стандартное уравнение} \quad \longrightarrow \quad
\text{отбор корней при необходимости}.
\]

Если в уравнении есть логарифмы от тригонометрических выражений, сначала учитываются ограничения логарифма, а затем решается полученное тригонометрическое уравнение.

\section*{2. Формулы понижения степени}

\[
\cos^2 x=\frac{1+\cos2x}{2},
\qquad
\sin^2 x=\frac{1-\cos2x}{2}.
\]

Важно: после понижения степени угол увеличивается в два раза.

Пример:
\[
\cos^2\left(x+\frac{\pi}{4}\right)
=
\frac{1+\cos\left(2x+\frac{\pi}{2}\right)}{2}
=
\frac{1-\sin2x}{2}.
\]

\section*{3. Формулы приведения}

Типовые переходы:
\[
\cos\left(\frac{\pi}{2}-x\right)=\sin x,
\qquad
\sin\left(\frac{\pi}{2}-x\right)=\cos x,
\]
\[
\sin\left(\frac{3\pi}{2}+x\right)=-\cos x.
\]

Формулы приведения удобно использовать, когда угол содержит
\[
\frac{\pi}{2},\quad \pi,\quad \frac{3\pi}{2},\quad 2\pi.
\]

Если угол содержит, например, \(\frac{\pi}{3}\), чаще нужна формула суммы или разности.

\section*{4. Формулы суммы и разности}

\[
\cos(\alpha-\beta)=\cos\alpha\cos\beta+\sin\alpha\sin\beta,
\]
\[
\sin\alpha+\sin\beta
=
2\sin\frac{\alpha+\beta}{2}\cos\frac{\alpha-\beta}{2}.
\]

Пример:
\[
\cos\left(x-\frac{\pi}{3}\right)
=
\cos x\cos\frac{\pi}{3}+\sin x\sin\frac{\pi}{3}
=
\frac12\cos x+\frac{\sqrt3}{2}\sin x.
\]

\section*{5. Как убрать «неудобный» минус перед \(x\)}

Если в аргументе косинуса стоит минус перед \(x\), можно использовать чётность:
\[
\cos(-u)=\cos u.
\]

Например:
\[
\cos\left(\frac{\pi}{4}-2x\right)
=
\cos\left(-(2x-\frac{\pi}{4})\right)
=
\cos\left(2x-\frac{\pi}{4}\right).
\]

Для синуса ситуация другая:
\[
\sin(-u)=-\sin u.
\]
Поэтому минус не исчезает, а выходит перед синусом.

\section*{6. Однородные тригонометрические уравнения}

Уравнение вида
\[
a\sin x+b\cos x=0
\]
часто решается делением на \(\cos x\), если \(\cos x\ne0\):
\[
a\tg x+b=0.
\]

Перед делением важно проверить, не теряются ли корни при \(\cos x=0\). Если \(\cos x=0\) не является решением исходного уравнения, деление допустимо.

Пример:
\[
\sqrt3\sin x-\cos x=0
\quad \Longleftrightarrow \quad
\sqrt3\tg x-1=0
\quad \Longleftrightarrow \quad
\tg x=\frac{1}{\sqrt3}.
\]

\section*{7. Логарифмы с тригонометрией}

Для уравнения
\[
\log_a F(x)=\log_a G(x)
\]
получаем
\[
F(x)=G(x),
\]
но при этом нужно учитывать:
\[
F(x)>0,\qquad G(x)>0.
\]

Иногда достаточно оставить только одно ограничение. Если \(F(x)=G(x)\) и \(F(x)>0\), то \(G(x)>0\) выполнится автоматически.

\subsection*{Пример ограничения}

Если после перехода есть условие
\[
\cos x>0,
\]
то работаем на правой полуокружности.

\begin{center}
\begin{tikzpicture}[scale=1.2]
\draw[->] (-1.35,0)--(1.45,0) node[right] {$\cos x$};
\draw[->] (0,-1.35)--(0,1.45) node[above] {$\sin x$};
\draw[darkaccent,line width=0.9pt] (0,0) circle (1);
\draw[accent,line width=2pt] (90:1) arc (90:-90:1);
\filldraw[white,draw=darkaccent,line width=0.9pt] (0,1) circle (1.8pt);
\filldraw[white,draw=darkaccent,line width=0.9pt] (0,-1) circle (1.8pt);
\node[right] at (1.05,0.55) {$\cos x>0$};
\end{tikzpicture}
\end{center}

\section*{8. Замена в логарифмических уравнениях}

Если один и тот же логарифм повторяется несколько раз, удобно сделать замену:
\[
t=\log_3(2\cos x).
\]

Тогда, например,
\[
2\log_3^2(2\cos x)-5\log_3(2\cos x)+2=0
\]
переходит в квадратное уравнение:
\[
2t^2-5t+2=0.
\]

После решения обязательно возвращаемся к исходной переменной \(x\).

\section*{9. Типичные ошибки}

\begin{itemize}
\item Забыть, что при понижении степени угол удваивается.
\item Применить формулу приведения к углу, где нужна формула суммы или разности.
\item Потерять знак при вынесении минуса из аргумента синуса.
\item Делить на \(\cos x\), не проверив случай \(\cos x=0\).
\item Забыть ОДЗ логарифма.
\item Писать «решений нет» там, где корни у промежуточного уравнения есть, но они не подходят по ограничению.
\item Не выделять на окружности нужную область при строгом ограничении.
\end{itemize}

\newpage

\section*{Тренировочный блок}

Решите уравнения.

\begin{enumerate}
\item \(\cos^2x=\frac12\).
\item \(\cos2x+\sin2x=1\).
\item \(\cos\left(x-\frac{\pi}{3}\right)-\cos x=0\).
\item \(\sin\left(\frac{\pi}{2}-x\right)+\sin x=1\).
\item \(\cos2x+\sin2x=0\).
\item \(\cos\left(x-\frac{\pi}{3}\right)-2\cos x=0\).
\item \(\log_2(\cos x)=\log_2(\sin2x)\).
\item \(2\log_3^2(2\cos x)-5\log_3(2\cos x)+2=0\).
\item \(\log_2(\cos x+\sin2x+8)=3\).
\item \(\cos^2\left(x+\frac{\pi}{4}\right)+\sin^2x=\frac12\).
\end{enumerate}

\newpage

\section*{Ответы}

\renewcommand{\arraystretch}{1.35}
\begin{tabularx}{\linewidth}{@{}cX@{}}
\toprule
№ & Ответ \\
\midrule
1 & \(x=\frac{\pi}{4}+\frac{\pi k}{2},\ k\in\Z\) \\
2 & \(x=\pi k\) или \(x=\frac{\pi}{4}+\pi k,\ k\in\Z\) \\
3 & \(x=\frac{\pi}{6}+\pi k,\ k\in\Z\) \\
4 & \(x=2\pi k\) или \(x=\frac{\pi}{2}+2\pi k,\ k\in\Z\) \\
5 & \(x=-\frac{\pi}{8}+\frac{\pi k}{2},\ k\in\Z\) \\
6 & \(x=\frac{\pi}{3}+\pi k,\ k\in\Z\) \\
7 & \(x=\frac{\pi}{6}+2\pi k,\ k\in\Z\) \\
8 & \(x=\pm\frac{\pi}{6}+2\pi k,\ k\in\Z\) \\
9 & \(x=\frac{\pi}{2}+\pi k\) или \(x=-\frac{\pi}{6}+2\pi k\) или \(x=\frac{7\pi}{6}+2\pi k,\ k\in\Z\) \\
10 & \(x=\pi k\) или \(x=\frac{\pi}{4}+\pi k,\ k\in\Z\) \\
\bottomrule
\end{tabularx}

\end{document}
